\chapter{Requirements} \label{Requirements}
The requirements were developed following analysis to ensure clearer understanding of the domain. These requirements are divided into functional and non-functional categories. Functional requirements describe the specific features or capabilities of the system, while non-functional requirements focus on \acrfull{QA}, as defined by Bass \cite{bass2021software}. According to Bass, functional requirements are tied to objective goals the system must achieve, whereas quality attributes describe the system's characteristics or properties. The scope of the functional requirements is further refined using the MoSCoW prioritization framework to ensure the product meets the criteria for a \acrfull{MVP}.

\section{Requirements engineering process}\label{SE_Engineering_process}
The project follows the Requirements engineering phases presented by Pressman and Maxim\cite{pressman2019software}.

\textbf{Inception:} The Inception Phase involved a meeting between the authors and two representatives from \acrshort{DEPA}, who served as a domain expert and a senior developer. During this meeting, a preliminary understanding of the desired product outcomes and the potential availability of data for testing purposes was established.

\textbf{Elicitation} The Elicitation Phase included interviews with domain experts and the examination of online sources to gain a deeper understanding of the domain. This ensured that the requirements would align with \acrshort{DEPA}'s expectations and that the final product would meet their envisioned use case. 

\textbf{Elaboration} The elaboration phase will be utilised to go deeper into the must have requirements, found trough the use of the MoSCoW framework. This knowledge will lay the foundation for the design phase. 

\textbf{Negotiation} There is no Negotiation phase as the primary goal of the project is to demonstrate the feasibility of \acrshort{DEPA}'s concept and not to deliver a fully functional product. This does pose a certain risk that the \acrshort{MVP} that is developed is not adequate for the user, this risk is considered acceptable given the project's exploratory nature.

\textbf{Validation} A validation phase will be conducted after each implementation milestone to ensure that the requirements have been fulfilled. This requires the functional requirements to be quantifiable in order to demonstrate that they have been met.

\section{Asking the first questions}
\label{SE_FirstQuestions}
After the First interview the Authors began utilising the engineering phases in the Pressman and maxim book\cite{pressman2019software}, and later found the section about asking the right questions to the stake holders and found them to be useful in getting a better understanding of the product and would have utilised the knowledge in the first meeting. and when this knowledge was acquired the Authors lost their stakeholder in the \acrshort{DEPA}. So a couple of questions have been picked out and answered based on the knowledge gotten trough chapter \ref{Analysis}, \ref{Related Work} and the Inception meeting.


\textbf{Who requested the work?} It has been a problem for some time that the way \acrshort{DEPA} tests if farmers use to much pesticides so it is \acrshort{DEPA} who requested this work and more specifically our main stakeholder kurt \textcolor{red}{jeg er ikke sikker på hvad fuck jeg skal gøre med hans navn eller om han overhovede skal være i reporten via navn eller bare main stakeholder} 

\textbf{Who will use the solution?} This question was hard to answer without the stakeholder, but it is important to understand the user to understand how the software should be utilised and how the data it produces should be presented. The user will properly not be so IT strong person since it will properly be the person doing the test or the one planing the tests. 

\textbf{What is the economic benefit of the solution?} This question was an easy one, it will better utilise the resource that is the test persons. So the resource is not wasted on random searches 

\section{Functional Requirements}\label{SE_FunctionalReqs}
The requirements were created independently from the stakeholder due to the fact there was negotiation phase, as explained in Section \ref{SE_FirstQuestions}. Even without direct involvement of stakeholders, most of the requirements were derived from the stakeholder interview and case files. The rest were derived by the knowledgeable gotten from Section \ref{Analysis} and \ref{Related Work}. 
This section will now explain the most important requirements, as seen in Table \ref{mustReqTable}. The rest of the requirements are found in Appendix \ref{Ap_LongReqTable}.
To start getting data out of the satellite the data about potato fields needs to be found, and \acrshort{DEPA} has a public Gis library over all farm land in Denmark, this is where (FR1) comes from. 
The main take away from the stakeholder resources are (FR2.3) and (FR3.3) where firstly the data is gathered and then calculate a kind of Overuse score. The way to calculate this Overuse Score is something that is not know before data is collected, and hold up against ground truth. (FR3.3.2) is about using an statistical model to calculate the overuse score. To ensure the data is correct a sanity check is needed (FR3.2) this check will make sure that there is data on the pixel. 
(FR5) is derived from the authors own experience from with downloading and processing big data. 

\begin{table}[H]
    \centering
    \begin{tabular}{|p{1.2cm}|p{6cm}|p{6cm}|}
    \hline
    \textbf{ID} & \textbf{Name} & \textbf{Description}\\
    \hline
    FR1 & Use the DEPA shapefiles & Read and unpack the shapefile dataset\\
    \hline
    FR1.1 & Extract Fields & Extract fields with the correct crop type \\
    \hline
    FR1.1.2 & Insert into DB & Store extracted fields in a database for NDVI association\\
    \hline
    FR2 & Download remote sensing data & Handle download of remote sensing data from the Sentinel-2 satellites\\
    \hline
    FR2.2 & API Authentication & Authenticate to Sentinel APIs before making calls \\
    \hline
    FR2.3 & Download images & Download satellite images with required bands for specific dates\\
    \hline
    FR3 & Process Satellite images & Handle image analysis workflows  \\
    \hline
    FR3.1 & Calculate NDVI & Calculate NDVI per pixel from image data  \\
    \hline
    FR3.2 & Validate pixel data & Ensure there is valid data before calculating NDVI \\
    \hline
    FR3.3 & Overuse score & Use NDVI to infer possible pesticide overuse  \\
    \hline
    FR3.3.2 & Statistical model & Calculate possible pesticide overuse based on a Statistical model \\
    \hline
    FR5 & Delegate work & Implement a work queue for long tasks like downloading and processing \\
    \hline
    \end{tabular}
    \caption{A table of all Must have requirements from the table found at appendix \ref{Ap_LongReqTable}}
    \label{mustReqTable}
\end{table}


\section{Scoping}\label{SE_Scoping}
The scope of the project was based on the stakeholder’s highest priorities—specifically, what they emphasized as essential for the product to function as intended. The only exception is that the stakeholder wanted to implement an \acrshort{AI} to calculate the overuse score for each field. However, since the primary goal of this project was to explore whether such a score could be observed and generated, a statistical approach was chosen instead.  Additionally, training an AI requires substantial amount if ground truth data, and there was an uncertainty if \acrshort{DEPA} could ever provide such ground truth.

\textbf{Why use MoSCoW} \\
MoSCoW is a prioritization tool commonly used in most of SDU's software courses. It categorizes requirements into four priority levels: \textit{(i) Must have:} These are essential functionalities required to deliver a viable \acrshort{MVP}. \textit{(ii) should have:} These features are important but not critical; they should be addressed only after all Must-haves are implemented. \textit{(iii) Could have:} These are desirable but non-essential features. They are not mission-critical and can be included if time and resources allow. \textit{(iv) Won't have this time:} These are out-of-scope features for the current iteration but may be added to the backlog for future consideration. 

\section{Non-Functional Requirements}\label{SE_NonFunctionalReqs}
The non functional requirements are described with \acrshort{QA} and \acrshort{QAS} to make them testable and more understandable to a non technical stakeholder\cite{bass2021software}.
The two \acrshort{QA} chosen are Scalability and Maintainability. Scalability was selected based on the stakeholders case that the collection of green index data is not limited to only on potatoes, but can be extended to other crop types as well. Maintainability was chosen by the authors to support the exploration aspect of the project. If the code is easy to modify with minimal errors, it significantly simplifies the process of exploring the data. \\
The structure of a \acrshort{QAS} consist of the following elements: \textit{(i) Source:} The person or system triggers the event in the system. \textit{(ii) Stimulus:} the event that occurs. \textit{(iii) Environment:} The state or condition the system is in when the event takes place. \textit{(iv) Artifact:} The part of the system that is effected. \textit{(v) Response:} The reaction to the event. \textit{(vi) Response Measure:} A way to evaluate the response, which could include metrics like response time or whether the system encounters an error.

\subsection{Scalability}
The main part of the scalability that is important to the project is to be able to change it from a small sample of potato field to a list of multiple crop types. 
So scalability is going to be measured in how well the system handles going from a single crop type to multiple as seen in \ref{tab:QASScala1}. Instead of increasing the data volume per API request, which would lead to larger and more complex responses. The work order is doubled to keep the API request more predictable in processing time and similar in token usage. 


\begin{table}[H]
    \centering
    \begin{tabular}{p{0.25\textwidth}|p{0.6\textwidth}}
         \hline
         \textbf{Element}& \textbf{Description}  \\
         \hline
         \textbf{Source}& A user who wants to test ndvi over potatoes compared to sugar beets\\
         \textbf{Stimulus} & Adding both potato and sugar beets to order system \\
         \textbf{Artifact} & Download and process system \\
         \textbf{Environment} & Normal mode. \\
         \textbf{Response} & It should generate work orders for both potato and sugar beets\\
         \textbf{Response Measure} & The work orders have doubled
    \end{tabular}
    \caption{QAS (1) for Scalability}
    \label{tab:QASScala1}
\end{table}

Scalability is also evaluated based on how easily the system can spin up new workers to better delegate work, as shown in \ref{tab:QASScala2}. This ensures that as demand increases, the system can be manually scaled up to handle the demand.
\begin{table}[H]
    \centering
    \begin{tabular}{p{0.25\textwidth}|p{0.6\textwidth}}
         \hline
         \textbf{Element}& \textbf{Description}  \\
         \hline
         \textbf{Source}& A user notices that the number of work orders is very high. Since sufficient computing resources are available, the user adds more workers to speed up the processing. \\
         \textbf{Stimulus} & Adding more download or process workers\\
         \textbf{Artifact} & Download and process system \\
         \textbf{Environment} & Normal mode. \\
         \textbf{Response} & New workers should spin up and start consuming work orders and logging to specific log file\\
         \textbf{Response Measure} & New workers become active and begin processing tasks within 30 seconds, with logs confirming startup and assignment of work. System throughput increases proportionally to the number of added workers, without errors or significant delays.
    \end{tabular}
    \caption{QAS (2) for Scalability}
    \label{tab:QASScala2}
\end{table}

\subsection{Maintainability}
Maintainability is evaluated based on how easy is it to add a new feature. The example would be adding a new green index to the code base as seen in Table \ref{tab:QAMain}

\begin{table}[H]
    \centering
    \begin{tabular}{p{0.3\textwidth}|p{0.6\textwidth}}
         \hline
         \textbf{Element} & \textbf{Description} \\
         \hline
         \textbf{Source} & A user wants to add support for EVI \\
         \textbf{Stimulus} & The data download logic must be updated to retrieve the required spectral bands, and a new calculation method must be added to the processing logic \\
         \textbf{Artifact} & The entire system \\
         \textbf{Environment} & Normal mode \\
         \textbf{Response} & A new push request is submitted with the necessary changes \\
         \textbf{Response Measure} & Integration of a new vegetation index should require no more than one day of development work
    \end{tabular}
    \caption{QAS for Maintainability}
    \label{tab:QAMain}
\end{table}



\section{Validation}\label{SE_Validation}
To validate the requirements, an accuracy evaluation will be performed using data provided by \acrshort{DEPA}. This evaluation aims to determine whether the system can accurately identify fields where pesticide usage exceeds acceptable thresholds. \\
Since the system relies on APIs and involves multiple interacting components, the overall functionality will be tested through integration testing. This testing will involve reviewing the logs of each module to ensure no errors occur during runtime, along with verifying that data is correctly recorded and displayed in the database.

